\documentclass[11pt]{article}

\usepackage[T1]{fontenc}
\usepackage{lmodern}
\usepackage{amsmath,amssymb,amsthm,mathtools}

\newtheorem{theorem}{Theorem}
\newtheorem{lemma}[theorem]{Lemma}
\theoremstyle{remark}
\newtheorem{remark}[theorem]{Remark}

\title{A Cubic Identity over Hyperfields}
\author{}
\date{}

\begin{document}
\maketitle

\begin{abstract}
We prove that if \(x,y,z\) are elements of a hyperfield \(\mathbb H\) and
\(0\in x\boxplus y\boxplus z\), then
\[
0\in x^3\boxplus y^3\boxplus z^3
\boxplus(-xyz)\boxplus(-xyz)\boxplus(-xyz).
\]
The proof uses the standard universal property of the near-regular partial field,
viewed as the pasture/tract presented by a three-term relation.
\end{abstract}

\section{Conventions and elementary facts}

Let \(\mathbb H\) be a hyperfield, with hyperaddition denoted by \(\boxplus\).
For a finite family \(a_1,\dots,a_n\in \mathbb H\), the expression
\(a_1\boxplus\cdots\boxplus a_n\) denotes the subset of \(\mathbb H\) obtained by
iterated hyperaddition; this is unambiguous by associativity of the underlying
canonical hypergroup.

We use the standard notation \(-a\) for the additive inverse of \(a\). Thus
\(-a=(-1)a\), and in particular \((-a)^3=-a^3\).

\begin{lemma}\label{lem:basic}
Let \(a,b,c\in \mathbb H\).
\begin{enumerate}
\item \(0\in a\boxplus b\boxplus c\) if and only if \(-c\in a\boxplus b\).
\item If \(0\in a_1\boxplus\cdots\boxplus a_n\) and \(\lambda\in \mathbb H\), then
\[
0\in \lambda a_1\boxplus\cdots\boxplus \lambda a_n .
\]
\end{enumerate}
\end{lemma}

\begin{proof}
For (1), by associativity,
\[
a\boxplus b\boxplus c=(a\boxplus b)\boxplus c.
\]
Thus \(0\in a\boxplus b\boxplus c\) holds precisely when there exists
\(w\in a\boxplus b\) such that \(0\in w\boxplus c\). By uniqueness of additive
inverses in a canonical hypergroup, this is equivalent to \(w=-c\).

For (2), multiply the inclusion \(0\in a_1\boxplus\cdots\boxplus a_n\) by
\(\lambda\) and use distributivity. Since \(\lambda\cdot 0=0\), the result follows.
\end{proof}

\section{A universal three-term relation}

We use the following standard presentation of the near-regular partial field.
Let
\[
\mathbb U
=
\mathbb F_{1}^{\pm}(A,B)\big/\!\!\big/\langle\!\langle A+B-1\rangle\!\rangle .
\]
Equivalently, \(\mathbb U\) is the partial field with ambient ring
\[
R=\mathbb Z[A,B,A^{-1},B^{-1}]/(A+B-1)
\]
and unit group generated by \(-1,A,B\). This is the usual near-regular partial
field. We use the standard fact that this presentation has the following universal
property: for every hyperfield \(K\), regarded as its associated tract, morphisms
\(\mathbb U\to K\) are naturally identified with pairs \((a,b)\in (K^\times)^2\)
such that
\[
0\in a\boxplus b\boxplus(-1),
\]
by sending \(A\mapsto a\) and \(B\mapsto b\). This is the standard description of
the foundation of \(U_{2,4}\); see Baker--Bowler for hyperfields as tracts
\cite{BB19} and Baker--Lorscheid for pasture presentations and the near-regular
partial field presentation \cite[Appendix A and Example 3.3]{BL24}.

\begin{lemma}\label{lem:nearregular}
Let \(K\) be a hyperfield, and let \(a,b\in K^\times\) satisfy
\[
0\in a\boxplus b\boxplus(-1).
\]
Then
\[
0\in
1\boxplus(-a^3)\boxplus(-b^3)
\boxplus(-ab)\boxplus(-ab)\boxplus(-ab).
\]
\end{lemma}

\begin{proof}
By the universal property above, there is a morphism of tracts
\[
\varphi:\mathbb U\longrightarrow K
\]
with \(\varphi(A)=a\) and \(\varphi(B)=b\). In the ambient ring of \(\mathbb U\) we
have \(A+B=1\), hence
\[
1-A^3-B^3-3AB=0.
\]
Indeed,
\[
(A+B)^3=A^3+B^3+3AB(A+B)=A^3+B^3+3AB.
\]
Thus the formal sum
\[
1+(-A^3)+(-B^3)+(-AB)+(-AB)+(-AB)
\]
is null in \(\mathbb U\). Applying \(\varphi\), and using that morphisms preserve
null relations, gives the asserted hyperfield inclusion.
\end{proof}

\begin{lemma}\label{lem:normalized}
Let \(K\) be a hyperfield, and let \(t,u\in K^\times\) satisfy
\[
u\in 1\boxplus t.
\]
Then
\[
0\in
1\boxplus t^3\boxplus(-u^3)
\boxplus tu\boxplus tu\boxplus tu.
\]
\end{lemma}

\begin{proof}
The inclusion \(u\in 1\boxplus t\) implies
\[
0\in 1\boxplus t\boxplus(-u).
\]
Multiplying this null relation by \(-1\), we obtain
\[
0\in (-1)\boxplus(-t)\boxplus u.
\]
Thus \(a=-t\) and \(b=u\) satisfy the hypothesis of Lemma~\ref{lem:nearregular}.
Applying that lemma gives
\[
0\in
1\boxplus(-(-t)^3)\boxplus(-u^3)
\boxplus(-(-t)u)\boxplus(-(-t)u)\boxplus(-(-t)u),
\]
which is exactly
\[
0\in
1\boxplus t^3\boxplus(-u^3)
\boxplus tu\boxplus tu\boxplus tu.
\]
\end{proof}

\section{The cubic identity}

\begin{theorem}
Let \(\mathbb H\) be a hyperfield, and let \(x,y,z\in \mathbb H\). If
\[
0\in x\boxplus y\boxplus z,
\]
then
\[
0\in
x^3\boxplus y^3\boxplus z^3
\boxplus(-xyz)\boxplus(-xyz)\boxplus(-xyz).
\]
\end{theorem}

\begin{proof}
First suppose that one of \(x,y,z\) is zero. By symmetry, assume \(x=0\). Then
\[
0\in 0\boxplus y\boxplus z=y\boxplus z,
\]
so \(z=-y\). Hence \(z^3=-y^3\) and \(xyz=0\). Therefore the desired hyper-sum
contains
\[
0\boxplus y^3\boxplus(-y^3)\boxplus 0\boxplus0\boxplus0,
\]
which contains \(0\). Thus the result holds whenever one of \(x,y,z\) is zero.

Now assume \(x,y,z\in\mathbb H^\times\). By Lemma~\ref{lem:basic},
\[
-z\in x\boxplus y.
\]
Set
\[
t=\frac{y}{x},
\qquad
u=\frac{-z}{x}.
\]
Multiplying \(-z\in x\boxplus y\) by \(x^{-1}\), we obtain
\[
u\in 1\boxplus t.
\]
Therefore Lemma~\ref{lem:normalized} gives
\[
0\in
1\boxplus t^3\boxplus(-u^3)
\boxplus tu\boxplus tu\boxplus tu.
\]
Multiplying this null relation by \(x^3\), Lemma~\ref{lem:basic} yields
\[
0\in
x^3
\boxplus x^3t^3
\boxplus x^3(-u^3)
\boxplus x^3tu
\boxplus x^3tu
\boxplus x^3tu.
\]
Now
\[
x^3t^3=x^3\left(\frac{y}{x}\right)^3=y^3,
\]
and, since \(u=-z/x\),
\[
x^3(-u^3)
=
-x^3\left(\frac{-z}{x}\right)^3
=
-(-z^3)
=
z^3.
\]
Finally,
\[
x^3tu
=
x^3\left(\frac{y}{x}\right)\left(\frac{-z}{x}\right)
=
-xyz.
\]
Substituting these identities gives
\[
0\in
x^3\boxplus y^3\boxplus z^3
\boxplus(-xyz)\boxplus(-xyz)\boxplus(-xyz),
\]
as required.
\end{proof}

\begin{thebibliography}{9}

\bibitem{BB19}
Matthew Baker and Nathan Bowler,
\emph{Matroids over hyperfields},
Advances in Mathematics \textbf{343} (2019), 821--863.

\bibitem{BL24}
Matthew Baker and Oliver Lorscheid,
\emph{On a theorem of Lafforgue},
International Mathematics Research Notices \textbf{2024}, no.~15 (2024), 11082--11091.

\end{thebibliography}

\end{document}